% =====================================================================
% 5. Model Validation and Result Analysis
% =====================================================================
\section{Model Validation and Result Analysis}
\label{sec:validation}

\begin{table}[htbp]
  \centering
  \caption{Summary of validation for the four problems}
  \label{tab:5-1}
  \small
  \begin{tabular}{@{}p{28mm}p{52mm}p{55mm}p{17mm}@{}}
    \toprule
    Object & Method & Deviation & Details \\
    \midrule
    Q1 global optimum & Closed-form $r^*$ vs.\ $0.1$~m grid & $7.5821$ vs.\ $7.5828$~m ($+0.0007$~m, $0.009\%$) & 4.2.4 \\
    Q2 accuracy and scheduling & Covariance vs.\ 200,000 MC; structural enumeration + top-5 + $0.5$~m variant & RMSE $0.1436$~m $=\sqrt2\,\mathrm{DRMS}$; P95 $0.2486$~m; optimum $10.355$~s, robust $10.837$~s & 4.3.2/4.3.5 \\
    Q3 error propagation & Linearization vs.\ $256\times256$ Hermite; $10^6$ MC & $s_{\max}$ $66.92$ vs.\ $66.84$~m ($0.13\%$); $P$ deviation $\le3\times10^{-4}$ & 4.4.3 \\
    Q4 rescheduling & 74-scenario accounting audit; $0.5$~m margin/tolerance/$22.5^\circ$ densification & 0 anomalies; $\eta/t_{\mathrm{last}}$ conclusions unchanged & 4.5.3 \\
    \bottomrule
  \end{tabular}
\end{table}

\subsection{Validation of Problem 1}

The closed-form lower bound $r^*=7.5821$~m differs from the $0.1$~m grid minimum $7.5828$~m by only $0.0007$~m ($0.009\%$); the observable area converges from $5610.4$ to $5608.5$~m$^2$ between the $0.1$~m and $0.05$~m midpoint grids ($0.03\%$ apart).

\subsection{Validation of Problem 2}

The covariance model gives $\mathrm{DRMS}=0.1016$~m; 200,000 MC runs give a radial RMSE of $0.1436$~m $=\sqrt2\,\mathrm{DRMS}$ and P95 $=0.2486$~m. The scheduling optimum $10.355$~s and the robust variant $10.837$~s corroborate each other.

\subsection{Validation of Problem 3}

The $P$ deviation of $10^6$ MC runs is $\le3\times10^{-4}$, and the linearization residual is $0.1293^\circ$ (about $0.2\%$ of the variance); the exact Hermite $s_{\max}=66.84$~m differs from the closed form $66.92$~m by $0.13\%$.

\subsection{Validation of Problem 4}

The 74-scenario accounting audit finds no anomalies; the $0.5$~m margin, zero-displacement tolerances of $0.02\bsl0.05\bsl0.10$~m, replacement with the $10.837$~s robust baseline, and $22.5^\circ$ densification all leave the $\eta/t_{\mathrm{last}}$ conclusions unchanged.

% =====================================================================
% 6. Sensitivity Analysis
% =====================================================================
\section{Sensitivity Analysis}
\label{sec:sensitivity}

\subsection{Parameter Selection and Perturbation Design}

Parameters are selected by relevance to the problem, dimensional scale, model dominance, and interpretability: geometric parameters $H$, $\varepsilon_{\max}$, $R_{\mathrm{veh}}$, and B2 height; accuracy parameters $\sigma_\alpha$, $d$, and $\gamma$; timeliness parameters $v$, $\omega$, and the 30~s limit; disturbance parameters $t_{\mathrm{rec}}$, failed vehicle, new-alert $\theta/\rho$, and safety margin. Continuous parameters use one-factor perturbations, $v\times\omega$ uses a $3\times3$ full factorial, discrete disturbances use the 74-scenario enumeration, and all other parameters stay at nominal values.

\subsection{Main Results}

\begin{table}[htbp]
  \centering
  \caption{Objective time under the $v\times\omega$ full factorial (unit: s)}
  \label{tab:6-1}
  \begin{tabular}{@{}lccc@{}}
    \toprule
    $v\backslash\omega$ & $24^\circ$/s & $30^\circ$/s & $36^\circ$/s \\
    \midrule
    $2.5$~m/s & 11.655 & 10.611 & 9.932 \\
    $3.0$~m/s & 11.499 & 10.355 & 9.676 \\
    $3.5$~m/s & 11.387 & 10.243 & 9.493 \\
    \bottomrule
  \end{tabular}
\end{table}

\noindent\textbf{Problem 1 elasticities.} The closed-form partial derivatives of $r^*$ are $\partial r^*/\partial H=0.1763$~m/m, $\partial r^*/\partial\varepsilon_{\max}=-0.7738$~m/deg, and $\partial r^*/\partial h_{\mathrm{cam}}=-0.1763$~m/m; at the nominal point $(H,\varepsilon_{\max})=(45~\mathrm{m},80^\circ)$, the elasticities are $E_H=1.05$ and $E_{\varepsilon\max}=(-0.7738\times80)/7.5821\approx-8.16$. As $H$ increases from 30 to 60~m, $r^*$ increases linearly from $4.94$ to $10.23$~m; as $\varepsilon_{\max}$ increases from $70^\circ$ to $85^\circ$, $r^*$ decreases nonlinearly from $15.65$ to $3.76$~m. Replacing $Q$ with $Q_2=(1/r)\cos^2(\varepsilon-45^\circ)$ leaves the argmax $(92.2,69.4)$, $r=7.5828$~m, and the optimal set unchanged; on the $\rmin$ circle, $Q_2$ is identically $0.0885$~m$^{-1}$, showing that the criterion choice is robust.

\noindent\textbf{Problem 2 elasticities.} The elasticity of $\omega$ is about $-0.44$ and that of $v$ about $-0.11$, so the former is about four times the latter; a $0.5$~m margin raises the objective from $10.355$ to $10.837$~s ($+4.7\%$) and increases the clearance from $1.3$~mm/$6.7$~cm to more than $0.5$~m. The sensitivity of $G^*$ to $H$ has the same origin as in Problem~1 ($0.1763$~m/m); $\sigma_\alpha$ and DRMS are strictly linear, and $G$ doubles at $\gamma=30^\circ$.

\noindent\textbf{Slope in $d$ for Problem 3.} $d=30\bsl35\bsl40\bsl45\bsl50\bsl60$~m corresponds to $s_{\max}=46.44\bsl56.68\bsl66.92\bsl77.16\bsl87.40\bsl107.88$~m, with a slope of $2.0481$~m/m; confidence levels $0.90\bsl0.95\bsl0.98\bsl0.99$ correspond to $85.20\bsl66.92\bsl51.41\bsl43.07$~m; gimbal-error scalings $0.5\bsl0.75\bsl1.0\bsl1.25\bsl1.5$ correspond to $72.32\bsl70.11\bsl66.92\bsl62.63\bsl57.03$~m. $d$ is the strongest continuous lever, the confidence level acts nonlinearly, and gimbal improvements saturate over long travels.

\subsection{Stability Conclusions}

Within reasonable perturbation ranges, no qualitative conclusion flips: the optimum lies on the boundary, $G^*=\rmin$, localization dominates over long travel, and rescheduling converts structural loss into temporal loss all remain valid. The model is insensitive to numerical/engineering parameters such as the grid step, $0.5$~m margin, zero-displacement tolerance, and $45^\circ$ azimuth sampling; the two-stage response of $t_{\mathrm{last}}$ to $t_{\mathrm{rec}}$ in Problem~4 is shown in Section~4.5.3. Remodeling is needed only under extreme conditions (the $\rmin$ circle is completely occluded, $t_{\mathrm{rec}}$ approaches 30~s with triple disturbances, or errors are strongly correlated).

% =====================================================================
% 7. Strengths and Weaknesses
% =====================================================================
\section{Strengths and Weaknesses}
\label{sec:strengths}

\subsection{Strengths}

\begin{enumerate}[label=(\arabic*)]
\item \textbf{Analytical certificates instead of numerical search.} Problem~1 derives $r^*=7.5821$~m from C2 and proves attainability on the feasible arc; Problem~2(1) proves $G^*=\rmin$ with the QM--AM inequality and C2; Problem~3 solves $s_{\max}$ in closed form from the monotonicity of $\sigma_\psi(s)$ and verifies it by exact quadrature. All three give a complete ``lower bound + attainability/equality'' proof rather than a best grid value.

\item \textbf{Integrated geometry--accuracy--timing--disturbance modeling with cross-validation.} The four problems share the same site geometry, observable region, and observation timing: $\Omega$ is used for station selection, the 2(2) timeline becomes the Problem~4 baseline, and the Problem~4 snapshot reuses the same slew/reset rules. Closed form--grid--sampling, covariance--MC--path simulation, linearization--Hermite--$10^6$ MC, and analytical benchmark--74 scenarios--robust re-solution all have residuals at the discretization/linearization level.

\item \textbf{Robustness and deployability built in.} Problem~2(2) discloses the $1.3$~mm/$6.7$~cm knife-edge clearance and gives a $10.837$~s robust variant; Problem~1 gives a $78^\circ$ fallback design with a $20.5\%$ cost; Problem~3 gives the conservative $66.84$~m value; Problem~4 uses 74 scenarios and $22.5^\circ$ densification for validation. Key numbers are driven by a single JSON data source, and the scripts are reproducible on any machine.
\end{enumerate}

\subsection{Weaknesses}

\begin{enumerate}[label=(\arabic*)]
\item \textbf{The global optimality of Problem~2(2) is not certified.} The assignment structures are exhaustively enumerated at the $10^3$ level, but station optimization relies on coordinate descent and Nelder--Mead; $10.355$~s is the best solution in the sense of ``structural exhaustion + station refinement,'' and the lower bound $3.414$~s is loose because it ignores reuse delays.

\item \textbf{The independent normal error assumption is strong.} Problem~2 assumes independent $\mathcal{N}(0,(1^\circ)^2)$ azimuth errors at the two stations, and Problem~3 assumes three independent zero-mean normal sources; MC validates the nominal parameters but does not cover correlated, nonzero-mean, or heavy-tailed cases.

\item \textbf{The scenario set is a discrete representative set, and the kinematics are simplified.} The 74 scenarios cover 8 azimuths, 3 distances, 5 failure types, and 3 recovery times, but not continuous arrivals, secondary disturbances, or adversarial choices; vehicles travel at constant speed with no acceleration/turning/collision avoidance, the gimbal rotates ideally, and the vehicle--target distance in Problem~3 is fixed at 40~m. These simplifications are consistent with the problem statement but would change the timeline and error propagation of a more realistic system.
\end{enumerate}

\subsection{Model Extensions}

\begin{enumerate}[label=(\arabic*)]
\item \textbf{Multi-station/multi-target.} The two-station covariance $\Sigma=A^{-1}DA^{-T}$ can be extended to multi-station intersection by stacking the rows of $A$; for multiple targets, each target has its own $\Sigma$ and shares the same vehicle and timeline constraints.

\item \textbf{Arbitrary monotone $Q(r)$.} The ``closed-form lower bound + attainability'' certificate of Problem~1 depends only on the monotonicity of $Q$ in $r$; the $Q_2$ comparison in Section~6.2 shows that switching to another monotone criterion in the same family leaves the optimal set unchanged, so the certificate transfers to other monotone quality measures.

\item \textbf{Continuous arrivals/secondary disturbances.} The event-triggered snapshot and re-optimization framework can directly accept continuously arriving new alerts and secondary disturbances; it suffices to refresh $A_p/V_p$ at the trigger time and reuse the same solver.
\end{enumerate}

% =====================================================================
% 8. Conclusion
% =====================================================================
\section{Conclusion}
\label{sec:conclusion}

First, the single-target observable region is determined by the intersection of the vehicle geometry, gimbal pitch, and LOS occlusion constraints; the optimal quality is fixed by the closed-form lower bound $r^*=7.5821$~m ($Q^*=0.1319$~m$^{-1}$), and the optimal set is a continuous arc that must be made unique by tie-breaking.

Second, two-station accuracy is characterized by the geometric dilution of precision $G$, whose lower bound is attained at $d_1=d_2=\rmin$ and $\gamma=90^\circ$; under the 30~s limit, V3+V5 is selected, and the maximum scheduling start is $10.355$~s, $50.6\%$ lower than the nearest-assignment baseline of $20.968$~s; timeliness and accuracy are not interchangeable.

Third, the pointing deviation is composed of localization, bias, and control sources; the maximum travel distance satisfying $P(|\Delta\psi|\le5^\circ)\ge0.95$ is $66.92$~m (with the exact-geometry conservative value $66.84$~m), and localization dominates over long travel.

Fourth, event-triggered rescheduling improves the completion rate in 67 of 74 scenarios, ties in 7, and degrades none; its value lies in converting structural loss into temporal loss.

Overall, geometry determines feasibility, accuracy determines station placement, scheduling determines timeliness, and rescheduling preserves the completion rate.

% =====================================================================
% References
% =====================================================================
\begin{thebibliography}{9}
\bibitem{swmcm} 2026 12th ShuWei Cup National College Mathematical Modeling Challenge (Autumn), Problem B: Rapid UAV Localization via Coordinated Mobile Vehicles. 2026.

\bibitem{preparata} Preparata F P, Shamos M I. Computational Geometry: An Introduction. New York: Springer, 1985.

\bibitem{torrieri} Torrieri D J. Statistical theory of passive location systems. IEEE Transactions on Aerospace and Electronic Systems, 1984, AES-20(2): 183--198.

\bibitem{foy} Foy W H. Position-location solutions by Taylor-series estimation. IEEE Transactions on Aerospace and Electronic Systems, 1976, AES-12(2): 187--194.

\bibitem{dijkstra} Dijkstra E W. A note on two problems in connexion with graphs. Numerische Mathematik, 1959, 1: 269--271.

\bibitem{nelder} Nelder J A, Mead R. A simplex method for function minimization. The Computer Journal, 1965, 7(4): 308--313.

\bibitem{hammersley} Hammersley J M, Handscomb D C. Monte Carlo Methods. London: Methuen, 1964.
\end{thebibliography}
